% ======================================================
%  Sekcja 3 --- Zagadki Logiczne
%  7 zadań | Sylogizmy, Pułapki słowne, Dedukcja, Proporcje
% ======================================================

% -- 19 --
\ExerciseSlide[\CategoryBadge[LogicColor!20]{Sylogizm}]
  {1}{2 min}{%
    Wszystkie Bloopy są Razzies.\par
    Wszystkie Razzies są Lazzies.\par
    Czy wszystkie Bloopy są Lazzies?
  }

\AnswerSlide[\CategoryBadge[LogicColor!20]{Sylogizm}]{Tak!}{
  {\footnotesize
    \textbf{Tak.}\\[0.4em]
    \begin{itemize}
      \item Bloopy $\subseteq$ Razzies
      \item Razzies $\subseteq$ Lazzies
      \item Stqd: Bloopy $\subseteq$ Lazzies \checkmark
    \end{itemize}
    To w\l{}asno\'s\'c przechodnio\'sci relacji podzbioru.
  }
}

% -- 20 --
\ExerciseSlide[\CategoryBadge[LogicColor!20]{Pu\l{}apka}]
  {2}{2 min}{%
    Rolnik ma 17 owiec.\par
    Wszystkie poza 9 zdychaj\k{a}.\par
    Ile owiec zostaje?
  }

\AnswerSlide[\CategoryBadge[LogicColor!20]{Pu\l{}apka}]{9 owiec}{
  {\footnotesize
    \textbf{Zostaje 9 owiec.}\\[0.4em]
    ``Wszystkie poza 9'' oznacza, \. ze 9 owiec \emph{nie} zdychaj\k{a}.
    Um.ys\l{} chce liczy\'{c}: $17 - 9 = 8$ --- to pu\l{}apka j\k{e}zykowa!
  }
}

% -- 21 --
\ExerciseSlide[\CategoryBadge[LogicColor!20]{Pu\l{}apka}]
  {2}{3 min}{%
    Dwie monety daj\k{a} razem 30 groszy.\par
    Jedna z nich NIE jest 10-groszowk\k{a}.\par
    Jakie to monety?
  }

\AnswerSlide[\CategoryBadge[LogicColor!20]{Pu\l{}apka}]{20~gr + 10~gr}{
  {\footnotesize
    \textbf{20-groszowka i 10-groszowka.}\\[0.4em]
    ``Jedna z nich nie jest 10-groszowk\k{a}'' --- prawda!\par
    20-groszowka \emph{nie} jest 10-groszowk\k{a}.\par
    Druga moneta jak najbardziej \emph{jest} 10-groszowk\k{a}.
  }
}

% -- 22 --
\ExerciseSlide[\CategoryBadge[LogicColor!20]{Dedukcja}]
  {2}{3 min}{%
    Adam, Bartek i Celina lubi\k{a} r\'{o}\. zne sporty:\par
    pi\l{}k\k{e} no\. zn\k{a}, p\l{}ywanie lub bieganie.\par
    Adam nie lubi pi\l{}ki. Bartek p\l{}ywa.\par
    Co lubi Celina?
  }

\AnswerSlide[\CategoryBadge[LogicColor!20]{Dedukcja}]{Pi\l{}k\k{e} no\. zn\k{a}}{
  {\footnotesize
    Krok po kroku:
    \begin{enumerate}
      \item Bartek $\Rightarrow$ p\l{}ywanie (podano wprost).
      \item Adam $\neq$ pi\l{}ka, Adam $\neq$ p\l{}ywanie $\Rightarrow$ Adam biega.
      \item Celina = jedyny sport bez przypisania $\Rightarrow$ \textbf{pi\l{}ka no\. zna}.
    \end{enumerate}
  }
}

% -- 23 --
\ExerciseSlide[\CategoryBadge[LogicColor!20]{Geometria}]
  {3}{5 min}{%
    Zegar pokazuje godzin\k{e} 3:15.\par
    Jaki k\k{a}t (w stopniach)\par
    tworz\k{a} wskaz\'{o}wki godzinowa i minutowa?
  }

\AnswerSlide[\CategoryBadge[LogicColor!20]{Geometria}]{$7{,}5^\circ$}{
  {\footnotesize
    \begin{itemize}
      \item Wskaz\'{o}wka minutowa o 3:15 $\Rightarrow$ $15 \times 6^\circ = 90^\circ$
      \item Wskaz\'{o}wka godzinowa o 3:00 $\Rightarrow$ $90^\circ$;
            w 15~min przesuwa si\k{e} o
            $\tfrac{15}{60} \times 30^\circ = 7{,}5^\circ$
            $\Rightarrow$ osi\k{a}ga $97{,}5^\circ$
      \item K\k{a}t: $97{,}5^\circ - 90^\circ = \mathbf{7{,}5^\circ}$
    \end{itemize}
  }
}

% -- 24 --
\ExerciseSlide[\CategoryBadge[LogicColor!20]{Proporcje}]
  {3}{5 min}{%
    5 maszyn robi 5 zabawek w 5 minut.\par
    Ile minut zajmie 100 maszynom\par
    zrobi\'{c} 100 zabawek?
  }

\AnswerSlide[\CategoryBadge[LogicColor!20]{Proporcje}]{5 minut}{
  {\footnotesize
    \textbf{5 minut.}\\[0.4em]
    Ka\. zda maszyna robi 1 zabawk\k{e} w 5 minut.\par
    100 maszyn r\'{o}wnolegle $\Rightarrow$ 100 zabawek w 5 minut.\\[0.3em]
    Pu\l{}apka: intuicja podpowiada $5 \times 20 = 100$ --- b\l{}\k{a}d!
  }
}

% -- 25 --
\ExerciseSlide[\CategoryBadge[LogicColor!20]{Pu\l{}apka}]
  {3}{5 min}{%
    W wy\'scigu wyprzedzasz osob\k{e}\par
    na 2.~miejscu.\par
    Na kt\'{o}rym miejscu teraz jeste\'s?
  }

\AnswerSlide[\CategoryBadge[LogicColor!20]{Pu\l{}apka}]{Na 2.~miejscu}{
  {\footnotesize
    \textbf{Na 2.~miejscu.}\\[0.4em]
    Wyprzedzaj\k{a}c kogo\'s na 2.~miejscu zajmujesz \emph{jego} pozycj\k{e}.\par
    Pierwszego miejsca nie tkn\k{a}\l{}e\'s!
  }
}
